\documentclass[12pt]{article}
\usepackage[utf8]{inputenc}
\usepackage[T1]{fontenc}
\usepackage[margin=1in]{geometry}
\usepackage{setspace}
\usepackage{xcolor}
\usepackage[colorlinks=true, urlcolor=blue!60!black, linkcolor=blue!60!black]{hyperref}
\setlength{\parindent}{0em}
\setlength{\parskip}{0.45em}
\renewcommand{\baselinestretch}{1.0}
\begin{document}
\begin{center}
{\Large \textit{Ex parte Young}, 209 U.S. 123 (1908)\par}
\vspace{0.3em}
{\normalsize Case Brief --- One-Page Edition\par}
\end{center}

\textbf{Memory Jogger.} Federal courts may enjoin state officers from enforcing unconstitutional laws; the Eleventh Amendment is no bar (\href{https://www.courtlistener.com/opinion/96819/ex-parte-young/}{CourtListener}).

\textbf{Facts.} Minnesota slashed intrastate railroad rates by commission order and statute, making passenger overcharges a felony and each violation of the 1907 Commodity Rate Act a misdemeanor punishable by up to 90 days in jail. Northern Pacific shareholders sued Attorney General Edward T. Young in federal court, alleging confiscatory rates and penalties that closed the courts.

\textbf{Procedural History.} Enjoined, Young sought state-court mandamus anyway; held in contempt and jailed, he sought habeas in the Supreme Court.

\textbf{Votes.} Petition dismissed 8--1. Peckham wrote for the Court; Harlan dissented.

\textbf{Holding.} A federal court may grant prospective relief against a state officer with some connection to enforcing an unconstitutional statute; the State cannot immunize him. Penalties that deter testing a law in court are facially unconstitutional.

\textbf{Reasoning.} Reaffirming \textit{Hans v. Louisiana}, Peckham nonetheless held that an officer enforcing an unconstitutional law acts without valid state authority and so has no lawful discretion to enforce it. Unlike the officers in \textit{Fitts v. McGhee}, Young's general duty to enforce state law supplied the needed connection, and awaiting prosecution was inadequate given the penalties. Harlan called the officer-State distinction a fiction that would ``practically obliterate the Eleventh Amendment'' and warned it would reduce States to ``dependencies''; state courts, bound by the same oath, could hear the railway's federal defenses. Peckham's remark that federal courts may not enjoin state courts themselves is influential dicta.

\textbf{Examples.} \emph{Same-Side:} enjoining a health commissioner's enforcement. \emph{Opposite-Side:} retroactive welfare payments from the treasury (\textit{Edelman}). \emph{Fence-Sitter:} privately enforced laws (\textit{Whole Woman's Health}).

\textbf{Critique.} The ``double fiction'' treats one act as state action for the Fourteenth Amendment but not the Eleventh, Progressives first decried ``government by injunction,'' yet \textit{Young} became the backbone of civil rights litigation.

\textbf{Key Quotation.} ``If the act which the state Attorney General seeks to enforce be a violation of the Federal Constitution, the officer, in proceeding under such enactment, comes into conflict with the superior authority of that Constitution, and he is, in that case, stripped of his official or representative character, and is subjected in his person to the consequences of his individual conduct.'' --- Justice Peckham, 209 U.S. at 159--60.

\end{document}
